\input format.tex

\usepackage{graphicx}
\graphicspath{{cores/}}

\begin{document}

\vspace*{6mm}
%% 各章节
\setlength{\arrayrulewidth}{1pt}
\fontsize{9.3pt}{11pt}\selectfont
\color{gray2}

\centerline{\bf\sanhao 健康建议}

\vspace*{2mm}

\begin{LRaside}[.20]{膳食方案}
{\bf *综合您肠道菌群检测结果，为您定制以下膳食方案：}\\
{\indent 建议清淡饮食，少食多餐，注意饮食卫生。以便秘为主者，可适当摄入燕麦、玉米等粗粮，增加富含膳食纤维的蔬菜水果，并足量饮水。腹泻为主者蔬菜水果可加工成蔬菜汤、果蔬汁等。蛋白类宜选用鱼虾、去皮禽肉等优质蛋白，多蒸煮，少煎炸。少油少盐，食用油控制在每天25克以内，食盐控制在每天6克以内。避免敏感食物，忌生冷、油腻及刺激性食物，如冷饮、酒精、浓茶、咖啡等，减少豆类、红薯、芋头、奶制品等产气食物的摄入。如患痛风、食物不耐受等有饮食禁忌的疾病，请综合考虑疾病的饮食原则。}\\
\asidebreak %
\noindent
\includegraphics[width=\linewidth]{shanshifangan.pdf}

\end{LRaside}


\begin{LRaside}[.70]{肠道调节方案}
\noindent
\includegraphics[width=\linewidth]{tiaojiefangan.pdf}

\asidebreak %
{\bf *综合您肠道菌群、营养功能检测结果，为您定制以下肠道调节方案：}\\
{\bf 益生元和益生菌}\\{\indent 1.每天服用适量低聚半乳糖或低聚木糖产品，促进体内双歧杆菌、乳酸杆菌等益生菌的生长，提高肠道菌群多样性。

2.每天服用适量益生菌产品（如酸奶、含活菌的益生菌饮料等），增加肠道内益生菌的数量，有利于肠道菌群平衡。

3.如患炎症性肠病、慢性腹泻，不宜服用菊粉及低聚果糖，建议选择低聚半乳糖。}\\
{\bf 补充辅酶Q}\\{\indent 帮助您增强抗氧化能力。}\\
{\bf 补充复合维生素}\\{\indent 均衡多种维生素营养，促进碳水化合物、蛋白质、脂类的正常代谢。}\\
\end{LRaside}

\begin{LRaside}[.20]{运动方案}
{\bf *适量运动可以帮助改善肠道菌群}\\
{\indent 可适当进行腹部搓揉，同时注意腹部保暖。不宜进行剧烈运动，可以根据自身恢复情况选择合适的运动方式，如散步、拉伸运动、羽毛球等。每次运动20-30分钟，每周3-5次。}
\asidebreak %
\noindent
\includegraphics[width=\linewidth]{yundongfangan.pdf}

\end{LRaside}

%\vspace*{0.5mm}
{\noindent\qihao *以上健康建议仅供参考，您的实际饮食、保健、运动等还需结合自身具体情况。}

\end{document}
